\section*{Part II: Analyze the Work Before Analyzing the Applicant}

\subsection*{11. Job Analysis Is the Intellectual Foundation of Selection}

\textbf{Job analysis} systematically identifies the activities, responsibilities, conditions, and competencies of work. It is distinct from writing an attractive job advertisement. An advertisement speaks to candidates. A job analysis tells the organization what it must predict. The U.S. Office of Personnel Management describes job analysis as the basis for linking work tasks, required competencies, and assessment procedures [R4]. This is general professional guidance, not a claim that U.S. public-sector rules govern an Indian startup.

At a small company, job analysis need not be a long bureaucratic exercise. The founder can inspect recent work, record critical incidents, identify recurring decisions, speak to people who have performed similar work, and estimate where errors arise. The output should be sufficiently specific that an outsider could recognize good performance. In larger organizations, job experts, task ratings, observation, and formal documentation become more important.

The underlying discipline is universal: \textbf{never create a test for a capability you cannot link to work the employee will actually do}. A beautiful assessment with no role connection is a distraction, however reliable its scores may be.

\subsection*{12. Start with Outputs and Critical Incidents}

A useful opening exercise is to reconstruct 10 representative work episodes: 5 unusually good and 5 unusually poor. An episode should describe the situation, the person's actions, the information available, the constraints, and the consequences. For an AI engineer, 1 incident might concern discovering data leakage before deployment. For an educator, it might concern a learner who appears confident but has a fundamental misconception. For a sales executive, it might concern a prospect pushing for an impossible contractual promise.

Critical incidents expose the choices that matter in real conditions. They also reveal whether the organization's supposed values are operationally coherent. If leadership says ``safety matters'' but routinely rewards shipping without testing, an applicant cannot infer what behavior will actually be valued. If academic rigor matters, the work should include refusing to pass inadequate work, documenting reasons, and supporting revision rather than merely keeping satisfaction ratings high.

The cases should sample mundane work too. Hiring exclusively for dramatic crises overlooks the ordinary consistency by which institutions earn trust: careful follow-up, timely feedback, accurate records, and repeated attention to detail.

\subsection*{13. Distinguish Tasks, Outputs, Competencies, and Prerequisites}

Consider a reviewer of technical student projects. \textbf{Task}: inspect code and experimental design. \textbf{Output}: a specific, accurate, actionable review delivered by deadline. \textbf{Competency}: recognizing model-evaluation errors and communicating corrections clearly. \textbf{Prerequisite}: enough statistics and software understanding to perform the review from day 1. These are connected but different things.

A job analysis should also distinguish \textit{entry} requirements from skills that can be trained. If the organization has a reliable 2-week procedure for teaching its grading system, prior familiarity with that exact rubric is not a justified screening criterion. If the job requires independent judgments about statistical validity immediately, that capability needs evidence before hiring.

A compact task-competency-assessment map illustrates this linkage:

\begin{table*}[t]
\centering
\small
\caption{Task-Competency-Assessment Map for Technical Review}
\label{tab:task_competency_map}
\begin{tabularx}{\textwidth}{p{3.5cm} p{3.8cm} p{3.8cm} X}
\toprule
\textbf{Important Work} & \textbf{Required Competency} & \textbf{Evidence} & \textbf{Risk If Missed} \\
\midrule
Identify flawed evaluation & Statistical and experimental judgment & Realistic flawed-model review & Learners receive false assurance \\
Give useful written feedback & Diagnostic communication & Marked-up learner submission & Learners cannot improve \\
Meet review deadlines & Planning and reliability & Work-history example, scoped trial & Support and trust degrade \\
Escalate uncertain cases & Integrity and calibrated judgment & Situational interview & Incorrect grades become entrenched \\
\bottomrule
\end{tabularx}
\end{table*}

Every row should have a logical link. If the evidence does not meaningfully measure the stated competency, revise it.

\subsection*{14. Rank Tasks by Importance, Frequency, Difficulty, and Consequences}

Frequent tasks are not necessarily the most consequential. A researcher may spend many hours cleaning data, but the ability to detect a fatal identification error may be more important than speed on routine cleaning. A customer-support agent may handle hundreds of easy questions, yet 1 rare mishandled privacy incident can be severe. A founder should therefore evaluate tasks along at least 4 dimensions: frequency, importance to outcomes, difficulty of competent performance, and consequences of failure.

Avoid creating an apparently precise formula by multiplying arbitrary 1-5 scores unless the result improves decisions. A simple criticality matrix may be enough: \textit{must perform independently on entry}; \textit{must perform after training}; \textit{nice to have}; \textit{not central}. Explicitly identify the non-negotiable tasks that cannot be compensated for by excellence elsewhere.

Task criticality also helps decide the assessment budget. The organization should spend more evaluation effort on competencies that are hard to observe after hiring, expensive to remediate, and closely tied to quality or safety. It need not examine every minor tool with equal intensity.

\subsection*{15. Write a Success Profile, Not a Wish List}

A weak specification reads like a catalogue of virtues: passionate, creative, hardworking, strategic, analytical, empathetic, proactive, detail-oriented, excellent communicator. Almost any applicant can claim these adjectives, and evaluators can interpret them however they wish. A \textbf{success profile} names the work that must be accomplished in a period, the competencies required, and the acceptable evidence of each.

For example: \textit{In the 1st 12 weeks, the selected reviewer can independently assess a defined category of project submissions, identify material reasoning errors, deliver specific feedback within agreed service levels, escalate ambiguous cases, and meet a documented calibration standard.} This describes a job the founder can assess. It can also guide onboarding and performance review.

Success profiles should be honest about uncertainty. If a new business has not decided the eventual volume or technology stack, state the likely scenarios and identify what can change. Do not promise a stable, narrowly defined job while quietly expecting the employee to rescue every new idea. Good candidates are evaluating the employer as carefully as the employer evaluates them.

\subsection*{16. Separate Minimum Standards from Differentiating Strengths}

A hiring rubric should not allow a candidate's exceptional public speaking to offset an inability to protect learner privacy. It may, however, allow excellent diagnostic teaching to offset merely adequate visual presentation. This suggests 2 classes of requirements. \textbf{Non-compensatory standards} must each be met: legal eligibility, necessary professional qualification, essential safety conduct, minimum technical correctness, relevant availability. \textbf{Compensatory strengths} can be combined: depth, speed, explanation style, creativity, process improvement, and collaboration, where these genuinely trade off in the role.

This distinction makes decisions more defensible. If a candidate fails a minimum standard, the reason is explicit. If 2 candidates both meet the threshold, the organization can compare their strengths relative to the work. It also prevents the post-hoc invention of a requirement when an interviewer dislikes someone.

Minimum standards should be few and justified. Excessive cutoffs reduce the pool and may systematically exclude talented people with different backgrounds. The founder must distinguish a true performance threshold from a preference for familiar credentials, aesthetics, or personality.
